% ===== §14 =====
\section{The Early Signifier and the Time of Relation}
\label{sec:timestructure}

\noindent
This section establishes the temporal structure that the persistence of drive gives to a relation. Its aim is to show that a bond is lived not only in its own present but under the pressure of signifiers inscribed long before it began, and that the two movements distinguished in Part Two, the circling drive and the sliding desire, together account for the peculiar shape of a relational history, at once faithful to a core it did not choose and carried forward into forms it could not have foreseen. The section sets out the early inscription, shows the two movements at work in the time of relation, and draws the claim about the shape of a relational history.

\subsection*{The early inscription}

Early in a life certain signifiers are inscribed with a weight the subject did not assign and cannot revoke. A name, a phrase, a scene given a form in the earliest field of the Other, these take on the office of organizing the subject's relation to the object-cause, and they do not lose that office when the subject grows and forms new bonds. This is the temporal reach of the account of Part Two. The remainder cut off in the constitution of the subject is not cut off once and then left behind; the manner of its cutting, the particular signifiers around which the subject's earliest relation to lack was organized, persists as a structure that later relations enter already shaped. Lacan's insistence that the unconscious is structured like a language has exactly this temporal consequence: a signifier laid down early does not fade with time but retains its structuring power, organizing what can be desired and how, long after its origin has been forgotten \citep{lacan2006ecrits}.

\subsection*{The two movements in the time of relation}

The later relation is entered by a subject already marked, already turned in a certain way around its void, and here the division of the two movements returns and does its work. The drive accounts for why the early core does not go out. Its circling is a repetition, and repetition is precisely the mode in which a thing persists without being renewed by choice; the early inscription endures because the drive keeps circling it, unspent across the years. The desire accounts for why the relation nonetheless keeps moving, why the subject does not simply repeat the same bond unchanged but slides through new figures, new partners, new forms of the same turning. Desire carries the subject forward along the chain, seeking in each new signifier what the last did not give, while the drive holds the center constant beneath the movement. A life of relation is the compound of the two: a constant core that the drive keeps circling and a displacement in time that desire keeps generating, the same manque holding beneath both.

\subsection*{The shape of a relational history}

This is what gives an intimate history its shape, neither pure repetition nor pure novelty. The point can be put as a claim, since it follows directly from the two offices established in Part Two. If it were only repetition, the drive alone, a person would form the same bond forever without development, condemned to reenact a single scene. If it were only displacement, desire alone, there would be no core, nothing that made a life of relation recognizably one person's, only a succession of unrelated attachments. The compound of a circling that persists and a sliding that displaces is what lets a relational life be at once continuous and in motion.

\begin{claim}[The shape of a relational history]
A relational history is neither pure repetition nor pure novelty but the compound of both. The drive's circling makes the early core persist, giving the history its continuity; desire's sliding carries the subject through new forms, giving the history its movement. A life of relation is recognizable as one person's because a core persists in it, and is a life rather than a fixed scene because it moves; the two movements together supply what neither could alone.
\end{claim}

\subsection*{Fidelity to the unchosen}

There is a consequence worth marking for what follows. The core a relational history is faithful to was not chosen; it was inscribed before choice, in the earliest field of the Other. The fidelity of a life to its core is therefore not, at bottom, a fidelity one elects, and this will bear on the ethics of the next section and of Part Four. To make room for another's core is to make room for something in them that they did not choose either, that was laid down in them before they could consent to it, and that organizes their desire from beneath. The time of relation is layered in this way for both partners, each entering the shared present under the pressure of a past inscription, and a relation is in part the meeting of two such layered times.
